% Comment conventions:
%   %% SCOPE:   what the section must cover (and must not)
%   %% FIGS:    planned figures/tables, with the script that makes them
%               (figures from ../output/figures/full/ or figures/; run the scripts with --dataset full)
%   %% NUMBERS: the script and output/results/full/*.json that every quoted number comes from
%   %% PAGES:   rough budget towards ~8 pages of body text (title page, references, appendices excluded)
%
% Compile: pdflatex -> biber -> pdflatex -> pdflatex   (Overleaf: pdfLaTeX; biber runs automatically).
% Missing figures compile as framed placeholders (see \plot below), so the draft builds before
% every script has been run on the full dataset.

\documentclass[a4paper,11pt]{article}
% ---------------------------------------------------------------- packages
\usepackage[margin=2.5cm]{geometry}
\usepackage[T1]{fontenc}
\usepackage[utf8]{inputenc}
\usepackage{lmodern}
\usepackage{microtype}
\usepackage{graphicx}
\usepackage{amsmath,amssymb}
\usepackage{bm}
\usepackage{siunitx}                    % \SI{9456.02}{\MeV}, \num{1.86e6}
\usepackage{booktabs}
\usepackage{array,tabularx}
\usepackage{caption}
\usepackage{subcaption}
\usepackage{xcolor}
\usepackage[hidelinks]{hyperref}
\usepackage[capitalise,noabbrev]{cleveref}
\usepackage[
  backend=biber,
  style=phys,
  articletitle=true,
  biblabel=brackets,
  chaptertitle=false,
  pageranges=false,
  sorting=none,
  maxbibnames=5, minbibnames=3,
  doi=true, eprint=true, url=true, isbn=false
]{biblatex}
\addbibresource{references.bib}

% Figures: straight from the scripts' output when compiled inside the repo; on Overleaf, upload
% output/figures/full/*.png into figures/.
\graphicspath{{figures/}{../output/figures/full/}}

% ---------------------------------------------------------------- float placement
\renewcommand{\topfraction}{0.85}
\renewcommand{\bottomfraction}{0.7}
\renewcommand{\textfraction}{0.15}
\renewcommand{\floatpagefraction}{0.8}
\setcounter{topnumber}{2}
\setcounter{bottomnumber}{2}
\setcounter{totalnumber}{4}

% ---------------------------------------------------------------- units and macros
\sisetup{separate-uncertainty = true, range-phrase = --, range-units = single, per-mode = symbol}
\newcommand{\MeVcc}{\ensuremath{\mathrm{MeV}/c^2}}
\newcommand{\GeVcc}{\ensuremath{\mathrm{GeV}/c^2}}
\newcommand{\GeVc}{\ensuremath{\mathrm{GeV}/c}}

\newcommand{\Ups}[1]{\ensuremath{\Upsilon(#1)}}               % \Ups{1S}
\newcommand{\mumu}{\ensuremath{\mu^+\mu^-}}
\newcommand{\pT}{\ensuremath{p_\mathrm{T}}}
\newcommand{\ptmin}{\ensuremath{\min(p_{\mathrm{T},1}, p_{\mathrm{T},2})}}
\newcommand{\NLL}{\ensuremath{-2\ln L}}
\newcommand{\iminuit}{\texttt{iminuit}}
\newcommand{\todo}[1]{\textcolor{red}{[TODO: #1]}}     % remove before submission

% \plot[width]{file}: the figure if it exists, otherwise a framed placeholder naming it.
\newcommand{\plot}[2][0.8\textwidth]{%
  \IfFileExists{figures/#2}{\includegraphics[width=#1]{#2}}{%
  \IfFileExists{../output/figures/full/#2}{\includegraphics[width=#1]{#2}}%
    {\fbox{\parbox[c][0.25\textheight][c]{#1}{\centering\texttt{\detokenize{#2}}\\(run the script on full data)}}}}}

% ---------------------------------------------------------------- title page
\begin{document}
\pagestyle{empty}
\pagenumbering{Alph}   % title page numbered A, so hyperref's page.1 anchor is not used twice

\IfFileExists{PandA_crest.pdf}{\par\noindent\includegraphics[width=12cm]{PandA_crest.pdf}}

\par\noindent
\vspace*{2cm}
\begin{center}
        \Large\bf Data Acquisition and Handling: Project F1\\[0.5em]
        \LARGE\bf Precision Measurement of the $\Upsilon(1S, 2S, 3S)$ Masses with LHCb Dimuon Data
\end{center}
\vspace*{0.5cm}
\begin{center}
        \bf Liam Kamphuis\\
\end{center}
\vspace*{5mm}

\begin{abstract}
%% SCOPE: 50–100 words. The aim (absolute Upsilon masses from extended maximum-likelihood fits),
%%        the method in one sentence (double Crystal Ball, selection, kinematic bins, systematics),
%%        the result (three masses with stat / fit-model / momentum-scale uncertainties) and the
%%        conclusion (momentum scale dominates; fit-model uncertainty below 1 MeV). Written LAST.
\todo{abstract}
\end{abstract}

\vfill
\todo{date}
\newpage

% ---------------------------------------------------------------- body
\pagestyle{plain}
\pagenumbering{arabic}
\tableofcontents
\newpage

% ================================================================ 1 INTRODUCTION
\section{Introduction}
\label{sec:intro}
%% SCOPE: the measurement (masses of the three Upsilon states from the 2015 LHCb dimuon sample)
%%        and why it is a test of a fitting procedure rather than of the masses themselves
%%        (PDG values are known to < 0.5 MeV); the aim in ONE sentence; what was done (one
%%        paragraph: Gaussian fits -> signal shapes from simulation -> selection -> kinematic bins
%%        -> systematics); main result in one or two sentences; roadmap.
%% NOT:   bottomonium spectroscopy beyond one sentence; LHCb detector beyond what Sec. 2 needs.
%% REFS:  LHCb:2013itw (the reference analysis), ParticleDataGroup:2024cfk.
%% PAGES: 0.5

The $\Upsilon$ mesons are bound states of a bottom quark and its antiquark ($b\bar b$).
The ground state, \Ups{1S}, was discovered in 1977 as a narrow peak in the dimuon mass
spectrum near \SI{9.5}{GeV/\clight^2}, the first evidence for a fifth quark. The radial
excitations \Ups{2S} and \Ups{3S} were found soon after. Because the $b$ quark is heavy,
the $b\bar b$ system is close to non-relativistic and is often called the hydrogen atom of
the strong interaction: the spacing of its levels probes the potential between a quark and
an antiquark. The three lowest states lie below the threshold for decay into a pair of $B$
mesons. Their widths are therefore tiny and they decay to $\mu^+\mu^-$ often enough to stand
out as three sharp peaks. Their masses are known to better than
\SI{0.5}{MeV/\clight^2}~\cite{ParticleDataGroup:2024cfk}. This makes them natural reference
points for a mass measurement: since the answer is known, what is tested is the analysis
itself, meaning how the peaks are modelled, how the uncertainties are estimated, and how
well the detector's momentum measurement is calibrated.



% ================================================================ 2 BACKGROUND
\section{Background}
\label{sec:background}
%% PAGES: 0.75 in total. Only what the method and the discussion refer back to.

\subsection{The \texorpdfstring{$\Upsilon$}{Upsilon} states and the LHCb dimuon sample}
\label{sec:bg-sample}
%% SCOPE: Upsilon(nS) as b-bbar bound states decaying to mu+mu-; natural widths (tens of keV)
%%        << detector resolution (tens of MeV), so the peak shape is the detector's, not a
%%        Breit-Wigner; the six observables per candidate (mass, pT, pseudorapidity eta,
%%        p, muon pTs); rapidity y vs the eta column; the background (combinatorial muon pairs).
%% REFS:  LHCb:2013itw; ParticleDataGroup:2024cfk.

The \Ups{1S}, \Ups{2S} and \Ups{3S} are the $n = 1, 2, 3$ radial states with the quark
spins aligned and no orbital angular momentum ($n\,{}^3S_1$). This gives them
$J^{PC} = 1^{--}$, the quantum numbers of the photon, so the $b\bar b$ pair can annihilate
through a virtual photon into a lepton pair. Decays into a $B\bar B$ pair are forbidden by
energy conservation below threshold (about \SI{10.56}{GeV/\clight^2}). The remaining strong
decays must go through three gluons, which is strongly suppressed (the OZI rule). As a
result the natural widths are only about 54, 32 and \SI{20}{keV}\todo{PDG 2024}, and the
dimuon branching fractions are 2--2.5\%\todo{PDG 2024}. At the LHC the $\Upsilon$ states
are produced mainly through gluon--gluon fusion. LHCb, a forward spectrometer covering
$2 < \eta < 5$, records them in large numbers~\cite{LHCb:2013itw}.


\subsection{Mass resolution, radiative tails and the momentum scale}
\label{sec:bg-resolution}
%% SCOPE: the mass is reconstructed from the muon momenta, so (i) the resolution follows
%%        delta m / m ~ delta p / p and depends on the muons' kinematics; (ii) final-state
%%        radiation gives a low-mass tail (Crystal Ball); (iii) a momentum-scale error
%%        p -> (1 + alpha) p shifts every mass by alpha m, the dominant systematic of LHCb mass
%%        measurements (calibrated to 3e-4).
%% REFS:  Skwarnicki:1986xj; LHCb:2023ood; LHCb:2011ikb (momentum-scale systematic in a mass
%%        measurement).

With natural widths of tens of keV, the observed width of each peak, tens of MeV, comes
almost entirely from the detector. The dimuon mass is computed from the two measured muon
momenta and their opening angle. At these momenta the momentum term dominates, so
$\delta m/m \approx \delta p/p$: the three peaks share the same relative resolution, and
their absolute widths grow with mass. The momentum itself is measured from how much each
track bends in the magnet. A faster track bends less, so its momentum is measured less
precisely, and the resolution varies with the muons' momenta and directions.

% ================================================================ 3 METHOD
\section{Method}
\label{sec:method}
%% PAGES: 1.5–2. Enough to make the results readable; detail goes to the appendices.

\subsection{Data, simulation and event selection}
\label{sec:m-data}
%% SCOPE: full 2015 sample (number of candidates, 9–11 GeV/c^2), simulated Upsilon(1S) sample and
%%        the generated masses (PDG 2016); the fiducial region of the paper (2 < y < 4.5,
%%        pT < 15 GeV/c) for the kinematic bins. The selection on the softer muon,
%%        min(pT1, pT2) >= 3 GeV/c, is motivated in Sec. 4.3; state here that tasks 1–4 use the
%%        inclusive sample and everything from the signal-shape study on uses the selected one.
%% NUMBERS: 01_measure_masses.py (event counts); 03_selection.py.
%% REFS:  ParticleDataGroup:2016lqr.

Our dataset comes from the LHCb experiment in 2015, with <> candidate dimuons spanning a mass range of 9 - 11 \GeVcc.
Our dataset is accompanied by a simulated Monte Carlo sample of the \Ups{1S} meson signal

\subsection{Extended maximum-likelihood fits}
\label{sec:m-fits}
%% SCOPE: extended likelihood (yields as parameters); binned (250 bins of 8 MeV) vs unbinned and
%%        why binned on data; every component normalised on the fit window so the composite PDF
%%        stays normalised; minimisation with MIGRAD, errors from HESSE, MINOS and -2 ln L scans;
%%        safeguards: convergence requires an accurate covariance and no parameter at a limit
%%        (false minimum of the double CB), restarts with SIMPLEX, starting values from a
%%        reference fit.
%% REFS:  Barlow:1990vc; James:1975dr; dembinski2025iminuit; Cowan:1998ji.

\todo{likelihood, minimisation, errors}

\subsection{Mass models}
\label{sec:m-models}
%% SCOPE: Gaussian (brief), double Gaussian, Crystal Ball and double Crystal Ball (baseline)
%%        signal; exponential background (quadratic as a systematic). The two constraints that
%%        make the three-peak fit identifiable: tails shared by all peaks; widths tied,
%%        sigma_k = sigma_1S m_k / m_1S (as in the reference analysis). Simultaneous fit in
%%        kinematic cells: masses and tails common, widths, yields and slopes per cell.
%% FIGS:  none; parameter definitions in a short table if space allows.
%% REFS:  Skwarnicki:1986xj; LHCb:2013itw.

\todo{signal and background shapes; shared and tied parameters}

\subsection{Goodness of fit, model comparison and validation}
\label{sec:m-gof}
%% SCOPE: Baker-Cousins chi2 for binned fits (and why not Pearson); runs test on the residual
%%        signs; likelihood-ratio test for nested models (Wilks, and its boundary caveat); AIC/BIC
%%        for non-nested ones; pseudo-experiments (toys) for bias and error coverage.
%% REFS:  Baker:1983tu; wald1940runs; wilks1938large; akaike1974new; schwarz1978estimating.

\todo{goodness of fit and validation}

\subsection{Systematic uncertainties}
\label{sec:m-syst}
%% SCOPE: shape bias measured on simulation and subtracted (each variation with its own bias);
%%        one-at-a-time variations (window, binning, signal and background shape, selection);
%%        why the shift of a cut variation is not judged with Barlow's formula (a cut removes
%%        mostly background, not a random subset) and the independent-slice test used instead;
%%        momentum scale 3e-4 as a separate, fully correlated uncertainty.
%% REFS:  Barlow:2002yb; LHCb:2023ood.

\todo{systematics}

% ================================================================ 4 RESULTS AND DISCUSSION
\section{Results and discussion}
\label{sec:results}
%% PAGES: ~5 with figures and tables. One subsection per study; each ends with what it means for
%%        the next. Every quoted number carries its uncertainty and comes from the named script.

\subsection{Single-Gaussian fits and errors by hand}
\label{sec:r-gauss}
%% SCOPE: brief tasks 1–4. 1S alone, then three peaks; parameter errors from HESSE, MINOS and
%%        -2 ln L scans, profiled vs others fixed, against sigma/sqrt(N); residuals and pulls
%%        showing the Gaussian fails (chi2/ndof ~ 13) with a structured pattern at each peak.
%% FIGS:  error_scans_1S_gauss.png; pulls_3peaks_gauss.png.
%% NUMBERS: 01_measure_masses.py --signal gauss --binned (fit_1S_gauss.json, fit_3peaks_gauss.json).

\begin{figure}[!htbp]
    \centering
    \plot[\textwidth]{error_scans_1S_gauss.png}
    \caption{\todo{caption: profiled vs fixed scans, HESSE band, Delta = 1}}
    \label{fig:error-scans}
\end{figure}

\begin{table}[!htbp]
    \centering
    \caption{Statistical errors on the \Ups{1S} Gaussian fit parameters [\MeVcc], by five methods:
    HESSE (curvature at the minimum), MINOS, the $\Delta(-2\ln L) = 1$ crossings of the profiled
    scan and of the scan with the other parameters fixed, and $\sigma/\sqrt{N}$
    ($\sigma/\sqrt{2N}$ for $\sigma_{1S}$), which ignores the background.}
    \label{tab:errors}
    \begin{tabular}{lccccc}
        \toprule
        & HESSE & MINOS & Profiled scan & Fixed scan & $\sigma/\sqrt{N}$ \\
        \midrule
        $\mu_{1S}$ & 0.100 & $^{+0.100}_{-0.100}$ & $^{+0.100}_{-0.100}$ & $^{+0.099}_{-0.099}$ & 0.059 \\
        $\sigma_{1S}$ & 0.103 & $^{+0.103}_{-0.103}$ & $^{+0.103}_{-0.103}$ & $^{+0.089}_{-0.089}$ & 0.042 \\
        \bottomrule
    \end{tabular}
\end{table}

\todo{text}

\subsection{Signal shape: simulation versus data}
\label{sec:r-shape}
%% SCOPE: brief task 5. MC (pure signal, known mass) shows the core and the FSR tail and each
%%        shape's bias; data peak wider than MC (more weight in the wide component); MC tails
%%        rejected by data (likelihood ratio); model comparison table (Gaussian -> double CB:
%%        chi2/ndof 12.9 -> 1.36); why the double CB is the baseline despite p(chi2) ~ 2e-4.
%% FIGS:  fit_3peaks_double_cb.png (with pulls); table: model comparison; table: MC vs data
%%        double-CB parameters.
%% NUMBERS: 02_signal_shape.py --binned (signal_shape.json).

\begin{table}[!htbp]
    \centering
    \caption{Three-peak fits to the selected data (250 bins) with each signal shape: free
    parameters $k$, Baker--Cousins $\chi^2$, its $p$-value, the $p$-value of the runs test on the
    signs of the residuals, AIC and BIC relative to the best model, and $\mu_{1S} - m_\mathrm{PDG}$
    [\MeVcc] (statistical error, no bias correction). Tails are shared by the three peaks; for the double CB
    with tails fixed to simulation, the likelihood ratio against free tails is
    $-2\Delta\ln L = 586$ for 4 parameters. The shape parameters are compared in \cref{tab:shape}.}
    \label{tab:comparison}
    \small
    \setlength{\tabcolsep}{4pt}
    \begin{tabular}{lccccccc}
        \toprule
        Signal & $k$ & $\chi^2/\mathrm{ndof}$ & $p(\chi^2)$ & $p(\mathrm{runs})$ & $\Delta$AIC & $\Delta$BIC & $\mu_{1S} - m_\mathrm{PDG}$ \\
        \midrule
        Gaussian & 11 & 3072/239 & $<10^{-300}$ & $3\times10^{-15}$ & 2745 & 2721 & $-4.47 \pm 0.09$ \\
        Double Gaussian & 11 & 1038/239 & $6\times10^{-100}$ & $3\times10^{-9}$ & 712 & 687 & $-4.42 \pm 0.09$ \\
        Crystal Ball & 11 & 971/239 & $7\times10^{-89}$ & $5\times10^{-9}$ & 645 & 620 & $-3.22 \pm 0.10$ \\
        Double CB, simulation tails & 9 & 908/241 & $5\times10^{-78}$ & $8\times10^{-7}$ & 578 & 528 & $-4.08 \pm 0.09$ \\
        Double CB & 13 & 322/237 & $2\times10^{-4}$ & $1\times10^{-3}$ & 0 & 0 & $-3.58 \pm 0.11$ \\
        \bottomrule
    \end{tabular}
\end{table}

\todo{text}

\subsection{Parameter correlations and over-parametrisation}
\label{sec:r-correlations}
%% SCOPE: the three degeneracies (width vs narrow fraction; alpha vs n; tail vs background) with
%%        their correlation coefficients, which appear in MC too and which only in data; the cost
%%        in parameter errors (sigma_1S, N_1S) vs the masses (barely affected: no mu in any strong
%%        correlation); how shared tails, tied widths and convergence checks handle them;
%%        profiled vs fixed scan width as a direct illustration.
%% FIGS:  possibly a correlation-matrix heat map of the baseline fit (not yet made).
%% NUMBERS: 02_signal_shape.py (correlation in signal_shape.json).

\todo{text}

\subsection{Event selection}
\label{sec:r-selection}
%% SCOPE: brief task 6. 2D plots; sideband estimate shows only the softer-muon pT separates
%%        signal from background (physics: Upsilon decays share momentum between both muons);
%%        cut scan: purity vs statistical mass error, flat optimum, round value 3.0 GeV/c;
%%        mass stability in independent slices (1S drifts ~3 MeV, chi2 80/6) as the
%%        motivation for the kinematic bins.
%% FIGS:  sideband_purity.png; cut_scan.png.
%% NUMBERS: 03_selection.py (selection.json).

\begin{figure}[!htbp]
    \centering
    \begin{subfigure}[b]{0.40\textwidth}
        \centering
        \plot[\linewidth]{sideband_purity.png}
        \caption{Sideband estimate.}
    \end{subfigure}\hfill
    \begin{subfigure}[b]{0.58\textwidth}
        \centering
        \plot[\linewidth]{cut_scan.png}
        \caption{Cut scan and mass stability.}
    \end{subfigure}
    \caption{\todo{caption}}
    \label{fig:selection}
\end{figure}

\todo{text}

\subsection{Fits in bins of transverse momentum and rapidity}
\label{sec:r-kinematic}
%% SCOPE: brief task 7. Grid (3 y x 4 pT) and what is shared; sigma_1S map (32–65 MeV, mostly
%%        y); masses vs inclusive (shift <= 0.12 MeV, errors 6–8% smaller); CB vs double-CB spread
%%        reduced; per-cell chi2 and the worst cell (curved background -> background-shape
%%        systematic).
%% FIGS:  kinematic_bins_sigma_double_cb.png; kinematic_bins_worst_cell_double_cb.png;
%%        table: masses inclusive vs simultaneous for both shapes.
%% NUMBERS: 04_kinematic_bins.py --signal double_cb / crystal_ball (kinematic_bins_*.json).

\begin{table}[!htbp]
    \centering
    \caption{Fitted masses minus PDG [\MeVcc] with statistical errors, for one inclusive fit and
    the simultaneous fit in 12 $(y, \pT)$ cells, in the fiducial region. $\chi^2$ is the
    Baker--Cousins statistic summed over all bins. No shape-bias correction is applied here.}
    \label{tab:kinematic}
    \begin{tabular}{llcccc}
        \toprule
        Signal & Fit & $\chi^2/\mathrm{ndof}$ & \Ups{1S} & \Ups{2S} & \Ups{3S} \\
        \midrule
        Crystal Ball & inclusive & 860/239 & $-3.22 \pm 0.11$ & $-3.68 \pm 0.27$ & $-4.00 \pm 0.47$ \\
        & 12 cells & 3666/2923 & $-3.38 \pm 0.10$ & $-3.86 \pm 0.25$ & $-4.25 \pm 0.44$ \\
        Double CB & inclusive & 312/237 & $-3.56 \pm 0.11$ & $-3.96 \pm 0.27$ & $-4.18 \pm 0.47$ \\
        & 12 cells & 3470/2921 & $-3.60 \pm 0.10$ & $-4.06 \pm 0.25$ & $-4.31 \pm 0.43$ \\
        \bottomrule
    \end{tabular}
\end{table}

\todo{text}

\subsection{Momentum scale}
\label{sec:r-scale}
%% SCOPE: data - MC peak position vs p (trend; a straight line does not describe it well) and vs
%%        softer-muon pT (flat: the slice drift is a resolution effect, not calibration); the
%%        overall offset (~ -4.5e-4) compared with the 3e-4 calibration uncertainty.
%% FIGS:  momentum_scale.png.
%% NUMBERS: 05_momentum_scale.py (momentum_scale.json).
%% REFS:  LHCb:2023ood.

\todo{text}

\subsection{Systematic uncertainties and final masses}
\label{sec:r-final}
%% SCOPE: brief task 8. Table of shifts per variation and their quadrature sum; final masses
%%        value +- stat +- fit model +- momentum scale and comparison with PDG; the 1S-calibrated
%%        2S and 3S masses and the common-scale fit as a cross-check; toy study: fit bias
%%        < 0.02 MeV, pull widths ~1 (errors trustworthy).
%% FIGS:  table: systematics; table: final masses; toy_pulls.png (or appendix).
%% NUMBERS: 06_final_masses.py (final_masses.json); 07_toys.py (toys.json).
%% REFS:  ParticleDataGroup:2024cfk.

\begin{table}[!htbp]
    \centering
    \caption{Uncertainties on the masses [\MeVcc]. Each fit-model row is the shift of the
    bias-corrected masses when one choice of the baseline fit is changed; a change of signal shape
    or selection is corrected with its own simulation fit.}
    \label{tab:systematics}
    \begin{tabular}{lccc}
        \toprule
        & \Ups{1S} & \Ups{2S} & \Ups{3S} \\
        \midrule
        Statistical (fit and simulation bias) & 0.13 & 0.27 & 0.44 \\
        \midrule
        Fit window 9.1--10.9\,GeV/$c^2$ & $+0.05$ & $+0.11$ & $+0.08$ \\
        Bin width doubled (16\,MeV) & $+0.02$ & $+0.03$ & $-0.07$ \\
        Crystal Ball signal & $-0.13$ & $-0.16$ & $-0.32$ \\
        Quadratic background & $-0.01$ & $+0.05$ & $+0.10$ \\
        $\ptmin \ge 2.5$\,GeV/$c$ & $+0.08$ & $-0.04$ & $-0.15$ \\
        $\ptmin \ge 3.5$\,GeV/$c$ & $+0.09$ & $-0.13$ & $-0.59$ \\
        No kinematic cells (inclusive fit) & $+0.07$ & $+0.14$ & $+0.17$ \\
        \midrule
        Fit model (sum in quadrature) & 0.20 & 0.28 & 0.73 \\
        Momentum scale ($3\times10^{-4}$) & 2.84 & 3.01 & 3.10 \\
        \bottomrule
    \end{tabular}
\end{table}

\begin{table}[!htbp]
    \centering
    \caption{Measured masses [\MeVcc]: value $\pm$ stat.\ $\pm$ fit model $\pm$ momentum scale, and
    the difference from PDG~\cite{ParticleDataGroup:2024cfk}. Below the line, the cross-checks of
    the momentum scale: the \Ups{2S} and \Ups{3S} masses with the scale calibrated on the \Ups{1S}
    (stat.\ only), and one scale factor fitted to all three peaks.}
    \label{tab:final}
    \begin{tabular}{lcc}
        \toprule
        & Mass & $-$ PDG \\
        \midrule
        \Ups{1S} & $9455.95 \pm 0.13 \pm 0.20 \pm 2.84$ & $-4.45$ \\
        \Ups{2S} & $10018.44 \pm 0.27 \pm 0.28 \pm 3.01$ & $-4.96$ \\
        \Ups{3S} & $10349.86 \pm 0.44 \pm 0.73 \pm 3.10$ & $-5.24$ \\
        \midrule
        \Ups{2S}, scale calibrated on \Ups{1S} & & $-0.25 \pm 0.27$ \\
        \Ups{3S}, scale calibrated on \Ups{1S} & & $-0.37 \pm 0.44$ \\
        Common scale $s-1$ & \multicolumn{2}{c}{$(-4.75 \pm 0.14)\times10^{-4}$, $\chi^2/\mathrm{ndof} = 1.44/2$} \\
        \bottomrule
    \end{tabular}
\end{table}

\todo{text}

% ================================================================ 5 CONCLUSION
\section{Conclusion}
\label{sec:conclusion}
%% SCOPE: the three masses and what limits them (momentum scale); what the fit studies showed
%%        (MC tails rejected, degeneracies confined to shape parameters, kinematic bins reduce the
%%        shape dependence); what would be done next (sigma parameterised in pT and y; a curved
%%        background per cell; a momentum-scale calibration of this sample). Half a page.

\todo{conclusion}

% ================================================================ AI USE
\section*{Declaration of AI use}
\addcontentsline{toc}{section}{Declaration of AI use}
%% SCOPE: what AI assistance was used for and what it was not, in line with the course guidance.
%%        Keep factual and short.

\todo{AI declaration}

% ================================================================ REFERENCES
\printbibliography[heading=bibintoc, title={References}]

% ================================================================ APPENDICES
\appendix

\section{Fit results}
\label{app:fits}
%% SCOPE: full parameter tables of the baseline fits (data and simulation) and the per-cell table
%%        of the kinematic fit; the toy pull figure if it does not fit in Sec. 4.
%% NUMBERS: signal_shape.json, kinematic_bins_double_cb.json, toys.json, bin_width_scan_gauss.json.

\begin{table}[!htbp]
    \centering
    \small
    \caption{The simultaneous double-CB fit per $(y, \pT)$ cell: selected candidates, fitted
    \Ups{1S} and \Ups{3S} yields, \Ups{1S} width, and the Baker--Cousins $\chi^2$ over the cell's
    250 bins (7--8 free parameters per cell plus the shared ones).}
    \label{tab:cells}
    \begin{tabular}{ccccccc}
        \toprule
        $y$ & \pT{} [GeV/$c$] & Candidates & $N_{1S}$ & $N_{3S}$ & $\sigma_{1S}$ [\MeVcc] & $\chi^2$ \\
        \midrule
        $2.0$--$3.0$ & $0$--$3$ & 335\,884 & 60\,812 & 5\,787 & $32.2 \pm 0.7$ & 273 \\
        $2.0$--$3.0$ & $3$--$6$ & 318\,742 & 66\,437 & 8\,854 & $33.0 \pm 0.7$ & 458 \\
        $2.0$--$3.0$ & $6$--$10$ & 138\,951 & 38\,235 & 5\,717 & $33.7 \pm 0.7$ & 309 \\
        $2.0$--$3.0$ & $10$--$15$ & 50\,383 & 18\,709 & 3\,499 & $35.8 \pm 0.8$ & 236 \\
        $3.0$--$3.5$ & $0$--$3$ & 225\,397 & 40\,375 & 4\,106 & $36.6 \pm 0.8$ & 242 \\
        $3.0$--$3.5$ & $3$--$6$ & 189\,646 & 40\,088 & 5\,381 & $37.1 \pm 0.8$ & 275 \\
        $3.0$--$3.5$ & $6$--$10$ & 110\,624 & 24\,586 & 3\,699 & $39.2 \pm 0.9$ & 249 \\
        $3.0$--$3.5$ & $10$--$15$ & 29\,649 & 11\,050 & 2\,172 & $41.9 \pm 0.9$ & 273 \\
        $3.5$--$4.5$ & $0$--$3$ & 163\,834 & 42\,671 & 4\,437 & $48.4 \pm 1.0$ & 299 \\
        $3.5$--$4.5$ & $3$--$6$ & 143\,537 & 43\,290 & 5\,496 & $50.2 \pm 1.0$ & 312 \\
        $3.5$--$4.5$ & $6$--$10$ & 72\,832 & 22\,749 & 3\,486 & $53.9 \pm 1.2$ & 269 \\
        $3.5$--$4.5$ & $10$--$15$ & 23\,925 & 10\,005 & 1\,951 & $64.4 \pm 1.5$ & 275 \\
        \bottomrule
    \end{tabular}
\end{table}

\begin{table}[!htbp]
    \centering
    \caption{Double-CB shape parameters on simulation (pure \Ups{1S}, unbinned), on data (one
    inclusive three-peak fit, tails shared) and in the simultaneous fit in cells. In the cell fit
    the narrow fraction is correlated with every cell's width ($+0.85$ to $+0.96$) and with the
    width ratio ($+0.92$), and $\alpha$ with $n$ ($-0.91$); no mass has a correlation above 0.5
    with any parameter.}
    \label{tab:shape}
    \begin{tabular}{lccc}
        \toprule
        & Simulation & Data, inclusive & Data, 12 cells \\
        \midrule
        $\sigma_{1S}$ [\MeVcc] & $34.3 \pm 0.1$ & $35.9 \pm 0.6$ & per cell (\cref{tab:cells}) \\
        $\alpha$ & $1.93 \pm 0.01$ & $1.52 \pm 0.04$ & $1.60 \pm 0.04$ \\
        $n$ & $1.08 \pm 0.02$ & $3.78 \pm 0.62$ & $3.66 \pm 0.58$ \\
        Width ratio $r$ & $2.16 \pm 0.03$ & $1.69 \pm 0.03$ & $1.52 \pm 0.06$ \\
        Narrow fraction $f$ & $0.907 \pm 0.005$ & $0.622 \pm 0.038$ & $0.730 \pm 0.068$ \\
        \bottomrule
    \end{tabular}
\end{table}

\begin{table}[!htbp]
    \centering
    \caption{Pseudo-experiments: 500 binned toys generated from the inclusive double-CB fit to
    the selected data and refitted (none failed). Pull $= (\hat\theta - \theta_\mathrm{true})/\hat\sigma$;
    the bias is the mean of $\hat\theta - \theta_\mathrm{true}$ [\MeVcc]; KS is the $p$-value of a
    Kolmogorov--Smirnov test of the pulls against a unit Gaussian.}
    \label{tab:toys}
    \begin{tabular}{lcccc}
        \toprule
        & Pull mean & Pull width & Bias & KS $p$ \\
        \midrule
        $\mu_{1S}$ & $-0.04 \pm 0.05$ & $1.014 \pm 0.032$ & $-0.005$ & 0.50 \\
        $\mu_{2S}$ & $-0.02 \pm 0.04$ & $0.990 \pm 0.031$ & $-0.004$ & 0.41 \\
        $\mu_{3S}$ & $+0.04 \pm 0.04$ & $0.984 \pm 0.031$ & $+0.017$ & 0.80 \\
        $\sigma_{1S}$ & $+0.10 \pm 0.04$ & $1.002 \pm 0.032$ & $+0.050$ & 0.23 \\
        \bottomrule
    \end{tabular}
\end{table}

\begin{table}[!htbp]
    \centering
    \caption{Bin width of the binned likelihood: the \Ups{1S} Gaussian fit (window 9.0--9.9\,GeV/$c^2$,
    $\sigma_{1S} = 43.3$\,\MeVcc) refitted with bins from 2 to 64\,MeV, compared with the unbinned fit
    (statistical error 0.0999\,\MeVcc). The three-peak fits use 8\,MeV bins.}
    \label{tab:bins}
    \begin{tabular}{ccccc}
        \toprule
        Bins & Width [MeV] & $\sigma_{1S}$/width & $\mu - \mu_\mathrm{unbinned}$ [\MeVcc] & Error / unbinned \\
        \midrule
        450 & 2.0 & 21.6 & $-0.0009$ & 1.0003 \\
        225 & 4.0 & 10.8 & $-0.0028$ & 1.0007 \\
        112 & 8.0 & 5.4 & $-0.0029$ & 1.0023 \\
        56 & 16.1 & 2.7 & $-0.0192$ & 1.0084 \\
        28 & 32.1 & 1.3 & $-0.0982$ & 1.0340 \\
        14 & 64.3 & 0.7 & $+0.0005$ & 1.1579 \\
        \bottomrule
    \end{tabular}
\end{table}

\section{Software and reproducibility}
\label{app:software}
%% SCOPE: half a page. The upsilon package (models, fitting, comparison) and the numbered scripts,
%%        one per study; Python and package versions (numpy, scipy, iminuit 2.33.0, particle
%%        1.1.1, mplhep); every figure and number regenerated with --dataset full.
%% REFS:  dembinski2025iminuit; rodrigues2025particle.

\begin{table}[!htbp]
    \centering
    \small
    \caption{Scripts and the sections they produce.}
    \label{tab:scripts}
    \begin{tabularx}{\textwidth}{>{\ttfamily}l X}
        \toprule
        \normalfont Script & Produces \\
        \midrule
        01\_measure\_masses.py & Gaussian fits, error scans, residuals (\cref{sec:r-gauss}) \\
        02\_signal\_shape.py   & simulation and data shape fits, model comparison, correlations (\cref{sec:r-shape,sec:r-correlations}) \\
        03\_selection.py       & sideband purity, cut scan, slice stability (\cref{sec:r-selection}) \\
        04\_kinematic\_bins.py & simultaneous fit in $(\pT, y)$ cells (\cref{sec:r-kinematic}) \\
        05\_momentum\_scale.py & data $-$ simulation peak position (\cref{sec:r-scale}) \\
        06\_final\_masses.py   & systematics and final masses (\cref{sec:r-final}) \\
        07\_toys.py            & pseudo-experiment validation (\cref{sec:r-final}) \\
        \bottomrule
    \end{tabularx}
\end{table}

\end{document}
